\documentclass[a4paper,11pt,twocolumn]{article}

\usepackage{luatexja}
\usepackage{luatexja-fontspec}
\setmainjfont{IPAexMincho}
\setlength{\textwidth}{175mm}
\setlength{\oddsidemargin}{-8mm}
\setlength{\textheight}{255mm}
\setlength{\topmargin}{-18mm}
\setlength{\abovedisplayskip}{4pt}
\setlength{\belowdisplayskip}{4pt}


\usepackage{graphicx}
\usepackage{here}
\usepackage{url}
\usepackage{cite}
\usepackage{amsmath}
\usepackage{titlesec}

\titlespacing{\section}{0pt}{0.5ex}{0.5ex}
\titlespacing{\subsection}{0pt}{0.5ex}{0.3ex}

\title{%
BLE電波強度を用いた屋内位置推定手法の精度評価と改善に関する研究\\
\large
}

\author{%
近畿大学 情報学部 情報学科\\
学籍番号：2312110159\\
氏名：内山敦也
}

\date{}

\begin{document}
\maketitle

\section{序論}
近年、スマートフォンや IoT デバイスの普及に伴い、屋内環境における位置情報サービスの需要が高まっている。
屋外では GPS による位置推定が一般的である一方、建物内部では電波遮蔽により測位精度が大きく低下するため、代替手法として
Wi-Fi、BLE（Bluetooth Low Energy）、UWB などの無線技術を用いた屋内測位手法が注目されている。

その中でも BLE は省電力・低コスト・スマートフォン標準搭載という利点から、実運用に適した技術として広く研究が進められている。
しかし、BLE 電波強度（RSSI: Received Signal Strength Indicator）は外乱の影響を強く受け、瞬間的なノイズや反射により
値が大きく変動するという問題がある。そのため、RSSI を距離に変換して位置推定に利用する手法では、
外れ値によって数 m ～ 数十 m の誤差が生じ、安定した測位精度を確保することが困難である。

本研究の目的は以下の 2 点である。
\begin{itemize}
    \item BLE RSSI のノイズ特性を考慮した平滑化手法を設計・実装し、測位精度を改善すること
    \item 実測データに基づき、平滑化処理の有無による位置推定精度を比較評価すること
\end{itemize}

本稿では、BLE 測位システムの構築方法、平滑化アルゴリズム、評価実験の手法および結果について報告し、
今後の改善点についても議論する。

\section{研究内容：BLEを用いた屋内位置推定手法}
本研究では，BLE（Bluetooth Low Energy）の電波強度 RSSI（Received Signal Strength Indicator）を用いて
屋内環境におけるタグ位置を推定するシステムを構築し，RSSI の変動特性に対する平滑化処理および
多点測位アルゴリズムの有効性を評価する。

本節では，(1) RSSI を距離に変換する受信モデル，(2) 位置推定に用いる最小二乗法に基づく多点測位，
(3) ノイズ低減のためのロバスト平滑化手法について述べる。

\subsection{RSSI から距離への変換モデル}
BLE電波の受信強度は，送受信間の距離だけでなく遮蔽物・反射・人体吸収・デバイスの向きによって大きく変動する。
一般に RSSI と距離 $d$ の関係は対数距離パス損失モデルにより次式で表される \cite{shuang2017ipin}。

\begin{equation}
d = 10^{\frac{(RSSI_{0} - RSSI)}{10n}},
\label{eq:pathloss}
\end{equation}

ここで，$RSSI_0$ は 1 m 離れたときの基準 RSSI，$n$ はパス損失指数（環境に依存し $1.6 \sim 3.5$ の値をとる）である。
BLE タグの向きが変化すると，電波の指向性の影響により RSSI が 3--10 dB 程度変動することがあり，
式 (\ref{eq:pathloss}) より距離推定誤差が 2--3 倍に拡大することが知られている。
本研究では ESP32/M5 デバイスを複数設置し，各アンカーからの RSSI を用いてタグの距離を推定する。

\section{今後の計画}
研究スケジュールをTable\ref{tab:schedule_phase}に示す。

\vspace{-0.8em}
\begin{table}[H]
\centering
\caption{研究計画（フェーズ分割）}
\label{tab:schedule_phase}
\small
\scriptsize
\setlength{\tabcolsep}{3pt}
\begin{tabular}{|c|p{5.6cm}|}
\hline
期間（フェーズ） & 研究内容 \\ \hline
フェーズ1（B3 3月〜B4 4月） &
既存研究調査および要件整理を行い，測位精度・想定利用人数・評価方法を明確化する。
システム全体の設計を行い，必要機材の選定・購入を実施する。 \\ \hline
フェーズ2（B4 5月〜6月） &
BLE測位システムの実装を行う。
M5Stamp/M5Tough の BLE スキャン処理，RSSI取得，サーバ通信を実装し，
平滑化手法および多点測位アルゴリズムを組み込む。 \\ \hline
フェーズ3（B4 7月〜9月） &
予備実験および実証実験を実施する。
静止・移動状態での精度評価と，複数デバイス同時測位時の動作安定性を確認し，
結果に基づきパラメータ調整を行う。 \\ \hline
フェーズ4（B4 10月〜11月） &
実験結果の整理と定量評価を行う。
従来手法との比較を通じて提案手法の有効性と課題を考察する。 \\ \hline
フェーズ5（B4 12月〜2月） &
卒業論文の執筆・最終調整を行い，論文提出および発表準備を行う。 \\ \hline
\end{tabular}
\end{table}

\vspace{-1.0em}
\section{参考文献}
\bibliographystyle{unsrt}
\bibliography{refs}

\end{document}
